% ===========================================================================
%  chapter1.tex  --  Глава 1. Теоретическая часть
% ===========================================================================

\section{Теоретическая часть}
\label{sec:theory}

\subsection{Временные ряды и постановка задачи}
\label{subsec:ts-problem}

\subsubsection{Введение в анализ временных рядов}
\label{sssec:ts-intro}

% TODO: Определение временного ряда, стационарность (строгая/слабая),
%       автоковариационная функция, разностный оператор.

\subsubsection{Постановка задачи прогнозирования}
\label{sssec:problem-statement}

% TODO: Формализация: входное окно L, горизонт H, модель f,
%       задача минимизации функции ошибки.

\subsubsection{Метрики качества прогноза}
\label{sssec:metrics}

% TODO: MSE, RMSE, MAPE, sMAPE --- определения и свойства.

\subsection{Классические модели временных рядов}
\label{subsec:classical-models}

\subsubsection{Модели AR, MA, ARIMA, SARIMA}
\label{sssec:arima}

% TODO: Формулировка AR(p), MA(q), ARMA(p,q), ARIMA(p,d,q), SARIMA.
%       Операторная запись, условия стационарности/обратимости.

\subsubsection{ACF, PACF и выбор порядка модели}
\label{sssec:acf-pacf}

% TODO: Определения автокорреляционной и частной автокорреляционной функций,
%       их использование для подбора порядков ARIMA.

\subsubsection{ETS (Exponential Smoothing)}
\label{sssec:ets}

% TODO: Простое экспоненциальное сглаживание, метод Хольта, метод Хольта--Уинтерса,
%       аддитивная и мультипликативная сезонность, пространство состояний.

\subsubsection{Модель Prophet}
\label{sssec:prophet}

% TODO: Аддитивная декомпозиция Prophet: тренд (кусочно-линейный или логистический),
%       сезонность (ряд Фурье), эффекты праздников.

\subsection{Нейросетевые модели для временных рядов}
\label{subsec:neural-models}

\subsubsection{N-BEATS}
\label{sssec:nbeats}

% TODO: Архитектура N-BEATS: стек блоков, basis expansion,
%       интерпретируемая конфигурация (trend + seasonality blocks).

\subsubsection{DLinear}
\label{sssec:dlinear}

% TODO: Идея Zeng et al. (2023): MA-декомпозиция + два линейных слоя,
%       аргумент о достаточности простых проекций.

\subsubsection{Autoformer}
\label{sssec:autoformer}

% TODO: Progressive decomposition architecture, Auto-Correlation mechanism,
%       сравнение с классическим self-attention.

\subsubsection{TimesNet}
\label{sssec:timesnet}

% TODO: Преобразование 1D -> 2D, multi-period Inception-блоки,
%       адаптивный выбор периодов через FFT.

\subsubsection{Helformer}
\label{sssec:helformer}

% TODO: Self-Attention + LSTM, Holt-Winters декомпозиция,
%       итеративный многошаговый прогноз.

\subsubsection{FEDformer}
\label{sssec:fedformer}

% TODO: Frequency Enhanced Decomposed Transformer,
%       Fourier/Wavelet attention, seasonal-trend decomposition.

\subsection{Обучаемая декомпозиция: сравнительный обзор}
\label{subsec:decomposition-comparison}

% TODO: Систематизация подходов к декомпозиции:
%       фиксированная (ETS, STL) vs обучаемая (DLinear MA-kernel,
%       Autoformer progressive, N-BEATS basis expansion).
%       Таблица сравнения.
